\documentclass{leveldoc}
\DocUses{foundations}
\DocFingerprint{4bd4d27b837cfd50}

\begin{document}

\section{What the level asks}\label{sec:asks}

\leveldescription

One positive ion is launched along $+x$ from the left edge, on the mid-line
between two metal plates. The upper plate is held at a negative potential and
the lower plate at the opposite positive potential, so the field in the gap
points upward and a positive ion is pushed upward. The plates are finite, so
the field bulges out at their ends, and a charge placed nearby rearranges the
charge on the metal. The player may add up to three charges, of either sign,
from the listed magnitudes, inside the marked region. There are no magnets.

The ion must enter the detector. The shot has no acceptance window, no kinetic
window and no gate. A flight ends when the ion touches a charge or a plate,
leaves the world, or enters the detector
(\xsecref{foundations}{sec:board}). A plate's obstacle is its box plus a
contact shell of thickness $\val{contact}$. A fixed charge has radius
$\val{Rcharge}$, and the ion's radius is $0$.

The reference adds one positive charge below the launch, at the near corner of
the region. The plates alone lift the ion, and it still leaves through the
right edge of the world below the detector. The added charge supplies the rest
of the upward push, and the ion meets the detector's lower edge. This document
states the boundary-element problem for the plates, gives the uniform-gap field
as the analytic limit, and integrates the flight in the panel field by the
classical Runge--Kutta method.

\section{Setup}\label{sec:setup}

\setuptable

The speed of light is $c = \val{cLight}$. The ion has charge $q = \val{qPart}$,
mass $m = \val{mPart}$ and radius $0$, and it is launched with kinetic energy
$T_0 = \val{T0}$. The launch is at $(\val{xLaunch}, \val{yLaunch})$, straight
along $+x$. Then $\gamma_0 = \val{gamma0}$, $\abs{\vect{p}_0} = \val{p0}$ and
$v_0 = \val{v0} = \val{vOverC}\,c$ (\xeqref{foundations}{eq:launch}).

The upper plate occupies $\val{xTop0} \le x \le \val{xTop1}$,
$\val{yTop0} \le y \le \val{yTop1}$, at potential $V_\mathrm{top} = \val{vTop}$.
The lower plate occupies $\val{xBot0} \le x \le \val{xBot1}$,
$\val{yBot0} \le y \le \val{yBot1}$, at potential $V_\mathrm{bot} = \val{vBot}$.
Both stand on the plane of the slice, symmetric about $z = 0$, with the height
given in the level file. The gap between the inner faces is $g = \val{gap}$.
The reference charge is $Q = \val{QPlay}$ at $(\val{xPlay}, \val{yPlay})$.

The detector is $\val{xDetLo} \le x \le \val{xDetHi}$,
$\val{yDetLo} \le y \le \val{yDetHi}$. The world is $-1 \le x \le \val{xWorld}$,
$-1 \le y \le \val{yWorld}$. The reported time step is $h = \val{H}$. The
reported mesh uses rectangles of edge $\val{panelFine}$, $\val{nPanels}$ of
them on the upper half of the two plates.

\begin{figure}[ht]
\centering
\begin{tikzpicture}
\begin{axis}[axis equal image, width=0.92\linewidth,
  xmin=-1.5, xmax=32.5, ymin=-2, ymax=22.5,
  xlabel={$x$}, ylabel={$y$}, tick label style={font=\small}]
  \draw[black!40, dashed] (-1,-1) rectangle (\rawval{xWorld},\rawval{yWorld});
  \draw[black!25] (0,0) rectangle (30,20);
  \fill[green!50!black, fill opacity=0.18]
    (\rawval{xDetLo},\rawval{yDetLo}) rectangle (\rawval{xDetHi},\rawval{yDetHi});
  \draw[green!40!black]
    (\rawval{xDetLo},\rawval{yDetLo}) rectangle (\rawval{xDetHi},\rawval{yDetHi});
  \fill[black!20]
    (\rawval{xTop0},\rawval{yTop0}) rectangle (\rawval{xTop1},\rawval{yTop1});
  \draw[black!70]
    (\rawval{xTop0},\rawval{yTop0}) rectangle (\rawval{xTop1},\rawval{yTop1});
  \fill[black!20]
    (\rawval{xBot0},\rawval{yBot0}) rectangle (\rawval{xBot1},\rawval{yBot1});
  \draw[black!70]
    (\rawval{xBot0},\rawval{yBot0}) rectangle (\rawval{xBot1},\rawval{yBot1});
  \fill[red!70!black] (\rawval{xPlay},\rawval{yPlay}) circle[radius=\rawval{Rcharge}];
  \addplot[orange!80!black] table {pathbare.dat};
  \addplot[very thick, blue!60!black] table {pathref.dat};
  \draw[-{Stealth[length=5pt]}, thick] (\rawval{xLaunch},\rawval{yLaunch}) -- (\rawval{arrowX},\rawval{arrowY});
  \fill (\rawval{xLaunch},\rawval{yLaunch}) circle[radius=0.08];
\end{axis}
\end{tikzpicture}
\caption{The plates and the two flights. The gray rectangles are the plates and
the disc is the reference charge; the green box is the detector. Blue is the
reference flight: it climbs through the gap and meets the bottom of the
detector. Orange is the same launch with that charge removed: the plates lift
it, and it leaves through the right edge of the world, below the detector.
Both paths stay outside the plates.}
\label{fig:board}
\end{figure}

\section{Principles and model}\label{sec:model}

\begin{principles}
  \item The field of a point charge \xeqref{foundations}{eq:coulomb}, and the
    force on a non-polarizable particle \xeqref{foundations}{eq:sphereforce}.
    There is no magnetic field.
  \item The relativistic equation of motion \xeqref{foundations}{eq:eom}, the
    launch \xeqref{foundations}{eq:launch}, and energy in a static field
    \xeqref{foundations}{eq:energy}. Away from charge the potential is harmonic,
    $\nabla^2\phi = -4\pi\rho$ as stated with that Coulomb field. The motion
    stays in the plane (\xsecref{foundations}{sec:plane}).
  \item A metal plate is an equipotential. Its relaxation is instantaneous on
    the flight's time scale (\physref{2.7}). The surface charge is constant on
    each rectangle of \S\ref{sec:panel}, fixed by collocation
    \eqref{eq:collocation}.
  \item The board's rules (\xsecref{foundations}{sec:board}): obstacles first,
    then the world, then the detector. The classical Runge--Kutta method and
    the cubic Hermite location of an event
    (\xsecref{foundations}{sec:numerics}).
  \item The radiated power \xeqref{foundations}{eq:larmor}, summed by Simpson's
    rule \xeqref{foundations}{eq:simpson}. It is not put back into the force.
\end{principles}

On the plates the potential equals the electrode potential. Off the metal the
electric field is the Coulomb field of the placed charges plus the field of the
surface charge and of its mirror in the plane $z = 0$. The ion is not part of
that surface charge. The state $(\vect{x}, \vect{p})$ obeys
\xeqref{foundations}{eq:eom} with $\vect{B} = \vect{0}$ and with that
$\vect{E}$, from the launch until the first ending event.

\begin{assumption}{Electrostatic plates}\label{as:static}
The metal answers at once. That drops the retardation across the board. The
launch speed is $\val{vOverC}\,c$. The game uses the same static plates
(\physref{2.7}), so this is the model being compared.
\end{assumption}

\begin{assumption}{No radiation reaction}\label{as:radiation}
The ion radiates \xeqref{foundations}{eq:larmor}, and that energy is not taken
from the flight. On the reference flight the radiated energy is a fraction
$\val{RefRadFrac}$ of $T_0$.
\end{assumption}

\begin{assumption}{No image of the ion}\label{as:image}
The ion's own surface charge on the plates is left out: a new solve at every
force evaluation. The game leaves it out for the same reason (\physref{2.7}).
The neglected force is bounded by $q^2/d^2$, with the distance $d$ to the metal
floored at $10^{-3}$, and the bound is taken relative to
$\max(\abs{q\vect{E}}, T_0)$. Along the reference samples the ratio reaches
$\val{RefImage}$; along the plates-alone samples it reaches $\val{BareImage}$.
The samples are the step ends and the ending point. The engine asks for the
ratio to stay below $10^{-10}$.
\end{assumption}

\begin{assumption}{Equal rectangles}\label{as:panels}
The surface charge is constant on equal rectangles of edge $\val{panelFine}$.
The engine collocates the same boundary condition on graded triangles
(\physref{2.7}). Replacing the edge by $\val{panelCoarse}$ moves an end state
by $\val{errMesh}$. That change is the discretization estimate of this mesh.
\end{assumption}

\begin{assumption}{The fixed step}\label{as:step}
The reported step is $h = \val{H}$. Halving it estimates an end state by
$\val{errStep}$ (\xsecref{foundations}{sec:numerics}). Ten times the panel-size
estimate $\val{errMesh}$ lies above $10^{-9}$, and it is larger than the
step-halving estimate $\val{errStep}$, so the tolerance on a state is that
product, $\val{tolFlight}$. Between samples the ion moves at most
$\val{stepReach}$, while the smallest sampled clearance of either plate is
$\val{RefDmin}$ on the reference flight and $\val{BareDmin}$ on the plates alone.
\end{assumption}

\section{Derivation}\label{sec:derivation}

\subsection{The uniform gap}\label{sec:gap}

Inside the gap, and away from the placed charges, $\nabla^2\phi = 0$. For two
infinite plates the potential depends only on $y$, so $\mathrm{d}^2\phi/\mathrm{d}y^2 = 0$ and
$\phi$ is linear in $y$. Matching the two potentials gives
$\mathrm{d}\phi/\mathrm{d}y = (V_\mathrm{top} - V_\mathrm{bot})/g$, and
$\vect{E} = -\nabla\phi$ then has
\begin{equation}\label{eq:uniform}
  E_y = \frac{V_\mathrm{bot} - V_\mathrm{top}}{g}.
\end{equation}
With these plates that value is $\val{eUniform}$, pointing from the positive
plate up toward the negative one. Equation \eqref{eq:uniform} is the field of
an infinite pair. These plates are finite: the panel field at the centre,
computed next, is lower.

With the plates alone, the panel field of \S\ref{sec:panel} is evaluated
halfway along the plates, on the launch line $y = \val{yLaunch}$, which is the
midpoint of the two centres. Its potential is $\val{phiCenter}$, the roundoff
of the cancellation between equal and opposite plates, and its upward field is
$\val{eCenter}$.
Relative to \eqref{eq:uniform} the difference is $\val{eUniformRel}$. The
flight is not integrated in the uniform field. The uniform field is the limit,
and $\val{eUniformRel}$ is how far the centre of this gap sits from that limit.

\subsection{One rectangle}\label{sec:panel}

A rectangle in the plane $z = z_0$, covering $x \in [x_0, x_1]$ and
$y \in [y_0, y_1]$, carries a constant density $\sigma$. Set
$u = x - x'$, $v = y - y'$, $w = z - z_0$ and $r = \sqrt{u^2 + v^2 + w^2}$.
Differentiating the logarithm gives the inner integral of $1/r$:
\begin{equation}\label{eq:log}
  \frac{\partial r}{\partial u} = \frac{u}{r}, \qquad
  \frac{\partial}{\partial u}\ln(u + r)
    = \frac{1 + u/r}{u + r} = \frac{1}{r}.
\end{equation}
The primitive in both variables, for the potential, is
\begin{equation}\label{eq:primitive}
  P(u, v, w) = u\ln(v + r) + v\ln(u + r)
    - \abs{w}\,\operatorname{atan2}\!\big(uv,\,\abs{w}\,r\big),
\end{equation}
where a logarithm whose argument is not positive is omitted. For $w > 0$ write
$\theta = \operatorname{atan2}(uv, wr)$, so the last term is $w\theta$. Its
derivative in $u$, and the two logarithms, are
\begin{align*}
  \frac{\partial}{\partial u}\big[u\ln(v + r)\big]
    &= \ln(v + r) + \frac{u^2}{r(v + r)}, \\
  \frac{\partial}{\partial u}\big[v\ln(u + r)\big]
    &= \frac{v}{r}, \\
  \frac{\partial}{\partial u}(w\theta)
    &= \frac{w^2 v}{r(u^2 + w^2)}.
\end{align*}
The angle derivative uses
$\partial\operatorname{atan2}(y, x)/\partial u = (x\,y_u - y\,x_u)/(x^2 + y^2)$
with $y = uv$ and $x = wr$: the numerator is
$(wv/r)(r^2 - u^2) = (wv/r)(v^2 + w^2)$, and the denominator is
$w^2 r^2 + u^2 v^2 = (u^2 + w^2)(v^2 + w^2)$. Collecting the three derivatives,
\begin{equation}\label{eq:p-u}
  \frac{\partial P}{\partial u}
    = \ln(v + r) + \frac{u^2}{r(v + r)} + \frac{v}{r}
      - \frac{w^2 v}{r(u^2 + w^2)}.
\end{equation}
The last three terms equal $u^2/(u^2 + w^2)$. Multiplied by $r$ they are
$u^2/(v + r) + v - w^2 v/(u^2 + w^2)$. The first two of those simplify by
$u^2 + v^2 = r^2 - w^2$:
\[
  \frac{u^2}{v + r} + v = r - \frac{w^2}{v + r},
\]
so the three terms, times $r$, equal
$r - w^2/(v + r) - w^2 v/(u^2 + w^2)$. This equals $r\,u^2/(u^2 + w^2)$ because
\[
  \frac{w^2}{v + r} + \frac{w^2 v}{u^2 + w^2} = \frac{r\,w^2}{u^2 + w^2},
\]
which clears to $u^2 + w^2 + v(v + r) = r(v + r)$, that is $r^2 + vr = r(r + v)$.
Therefore
\begin{equation}\label{eq:p-u-simple}
  \frac{\partial P}{\partial u} = \ln(v + r) + \frac{u^2}{u^2 + w^2}.
\end{equation}
The second term does not depend on $v$. The same reduction as \eqref{eq:log},
with $v$ in place of $u$, gives $\partial\ln(v + r)/\partial v = 1/r$, so
\begin{equation}\label{eq:mixed}
  \frac{\partial^2 P}{\partial v\,\partial u} = \frac{1}{r}.
\end{equation}
Integrating \eqref{eq:mixed} from $u = x - x_1$ to $u = x - x_0$ and from
$v = y - y_1$ to $v = y - y_0$ produces the corner sum below. Those limits are
the rectangle, because $\mathrm{d}x' = -\mathrm{d}u$ and $\mathrm{d}y' = -\mathrm{d}v$:
\begin{equation}\label{eq:rectangle}
\begin{aligned}
  \phi &= \sigma\Big[
      P(x - x_0, y - y_0, w) - P(x - x_0, y - y_1, w) \\
      &\quad - P(x - x_1, y - y_0, w) + P(x - x_1, y - y_1, w)
    \Big].
\end{aligned}
\end{equation}
The field is $\vect{E} = -\nabla\phi$, from the same corner sum. Replacing $w$
by $\abs{w}$ in the last term of \eqref{eq:primitive} makes $P$ even in $w$ and
leaves \eqref{eq:mixed} unchanged for $w \neq 0$. The potential of a sheet is
therefore the same on both sides, the in-plane field is even, and the normal
field is odd. On the sheet itself $1/r$ is integrable, so the potential at a
centroid, which lies in the interior of its rectangle, is finite.

\subsection{The plates}\label{sec:plates}

Each plate is a box. Only the upper half is meshed: the four side faces from
$z = 0$ up to half the height, and the top face. The mid-plane is the interior
of the metal, so it is not a face. Each edge from $a$ to $b$ is split into
$n = \max(1, \operatorname{round}((b - a)/\val{panelFine}))$ equal parts, and
the face is the grid of those parts. The mirror of a panel through $z = 0$
carries the same density and is not an extra unknown. The reported mesh has
$\val{nPanels}$ upper rectangles.

Let $\vect{r}_i$ be the centroid of upper rectangle $i$, belonging to electrode
$e(i)$, and let $\phi_j$ be the potential \eqref{eq:rectangle} of that
rectangle at unit density. Collocation asks
\begin{equation}\label{eq:collocation}
  \sum_j \sigma_j \big(\phi_j(\vect{r}_i) + \phi_j^\mathrm{mirror}(\vect{r}_i)\big)
    = V_{e(i)} - \phi_\mathrm{ext}(\vect{r}_i).
\end{equation}
Here $\phi_\mathrm{ext}$ is the Coulomb potential of the placed charges. The
ion is absent from $\phi_\mathrm{ext}$ and from the panels, which is
assumption~\ref{as:image}. A placed charge changes every density through the
right-hand side: the plates' charge rearranges. With no placed charge,
$\phi_\mathrm{ext} = 0$. The densities are the solution of this dense linear
system. The residual $\max\abs{G\sigma - b}$ on the reported mesh is
$\val{solveRes}$.

Off the centroids the piecewise-constant density does not hold the potential
exactly. On the inner face of each plate, in the plane $z = 0$, at two stations
that miss the centroids, and with the reference charge present, the potential
differs from the plate potential by at most $\val{surfRes}$.

In the plane $z = 0$ the normal fields of a panel and its mirror cancel, so
$E_z = 0$ and the ion stays in the plane. The field on the ion is the in-plane
field of the panels and the mirrors, plus the Coulomb field of the placed
charges.

\subsection{The flight}\label{sec:flight}

There is no closed orbit. The state is advanced by \xeqref{foundations}{eq:rk4}
with the fixed step $h = \val{H}$. After each step the cubic Hermite
interpolant of the two ends is sampled at thirty-two fractions of the step. The
first sample that has touched an obstacle, left the world or entered the
detector is refined by bisection, sixty halvings
(\xsecref{foundations}{sec:numerics}). Obstacles are tested before the world,
and the world before the detector. A face coordinate is assigned when the
residual is below $10^{-6}$.

The conserved energy is \xeqref{foundations}{eq:energy}. The rest energy
$mc^2$ does not depend on the state, so the diagnostic is the kinetic energy
plus the potential of the panels, the mirrors and the placed charges,
\begin{equation}\label{eq:flight-energy}
  W = (\gamma - 1)mc^2 + q\,\phi(\vect{x}).
\end{equation}
There is no image term: the ion is not in $\phi$. The error of the fine step
is $\abs{y_h - y_{2h}}/15$. The same rule estimates the Simpson error of the
radiated energy. The panel-size estimate is the change in the end state when
the rectangle edge doubles, not divided by $15$. Both flights give the same
outcome at $h$ and at $2h$, and on both meshes.

\section{The reference solution by hand}\label{sec:hand}

\begin{procedure}
  \item\label{st:launch} \emph{The launch} \xeqref{foundations}{eq:launch}.
    $\gamma_0 = \val{gamma0}$, $\abs{\vect{p}_0} = \val{p0}$, so
    $\vect{p}_0 = (\val{p0}, 0)$ and $v_0 = \val{v0}$.
  \item\label{st:gap} \emph{The gap, as a limit.} Equation \eqref{eq:uniform}
    gives $E_y = \val{eUniform}$. The panel field at the centre of the gap,
    with the plates alone, is $\val{eCenter}$, lower by the relative amount
    $\val{eUniformRel}$. The potential there is $\val{phiCenter}$. The march
    uses the panels, not this uniform value.
  \item\label{st:panels} \emph{The panels.} Build the upper-half rectangles of
    edge $\val{panelFine}$ as in \S\ref{sec:plates}: $\val{nPanels}$ of them,
    plus the mirrors. Solve \eqref{eq:collocation} with the reference charge on
    the right-hand side. The residual at the centroids is $\val{solveRes}$. On
    the inner faces, off the centroids, the potential is off by at most
    $\val{surfRes}$.
  \item\label{st:march} \emph{The march.} Advance \eqref{eq:flight-energy} and
    \xeqref{foundations}{eq:eom} by \xeqref{foundations}{eq:rk4} with step
    $\val{H}$, and locate the ending face as in \S\ref{sec:flight}. The state
    every five time units:
    \begin{center}\small\val{states}\end{center}
  \item\label{st:end} \emph{The arrival.} The ion meets the \val{RefFace} face
    of the detector at $(\val{RefX}, \val{RefY})$, at $t = \val{RefT}$, with
    momentum $(\val{RefPx}, \val{RefPy})$, at an angle $\val{RefAng}$ degrees
    above $+x$. The ordinate is the detector's lower edge. The run takes
    $\val{RefSteps}$ steps. The path only climbs: its lowest point is the
    launch, and its highest is $\val{RefYmax}$.
  \item\label{st:check} \emph{The checks.} Along the flight, $W$ drifts by at
    most $\val{RefWdev}$. The step-halving estimate of the end state is
    $\val{RefErr}$, and the larger of the two flights is $\val{errStep}$. The
    panel-size estimate is $\val{errMesh}$. The sampled clearance of the plates
    is $\val{RefDmin}$, outside the shell $\val{contact}$, and the ion moves at
    most $\val{stepReach}$ in one step. The image-force ratio reaches
    $\val{RefImage}$.
  \item\label{st:rad} \emph{The radiation.} Simpson's rule on steps of $h/2$
    gives $\val{RefRad}$, a fraction $\val{RefRadFrac}$ of $T_0$. Halving that
    step again changes the integral by the relative amount $\val{RefRadErr}$.
    This is our field's integral. It is not a row of the comparison: the
    engine's triangles are a different field.
\end{procedure}

\begin{result}[The reference flight]
The ion enters the detector through its lower edge at
$(\val{RefX}, \val{RefY})$ at $t = \val{RefT}$, with momentum
$(\val{RefPx}, \val{RefPy})$. The level is solved.
\end{result}

The march is a few thousand steps of four field evaluations. Each evaluation
sums the $\val{nPanels}$ rectangles, their mirrors, and the one placed charge.

\section{Comparison with the game}\label{sec:comparison}

\comparisontable

Eight comparisons are in the table: the reference outcome, time and end
position, and the outcome, time and end position of the plates with no added
charge. The tolerance on a state is $\val{tolFlight}$. Ten times the panel-size
estimate $\val{errMesh}$ lies above $10^{-9}$, and it is larger than the
step-halving estimate $\val{errStep}$, so the tolerance is that product. The
tolerance on a state was fixed from these estimates before the comparison.

The largest absolute difference is the bare ordinate, $\val{cmpBareYDiff}$. It
matches the panel-size estimate $\val{errMesh}$: the two are the same size, and
the difference lies under the tolerance. The next is the reference abscissa,
$\val{cmpRefXDiff}$, the same size again. The reference time differs by
$\val{cmpRefTDiff}$ and the bare time by $\val{cmpBareTDiff}$, both smaller.
The step estimate $\val{errStep}$ is far below all four: the panels, not the
time step, set the difference. The outcomes agree.

The reference ordinate and the bare abscissa agree exactly. Each is a face the
flight meets, the detector's lower edge and the world's right edge.

\section{Other solutions}\label{sec:others}

The plates alone, which is the empty placement in \code{placements.json}, do
not solve the level. There is no slider, and the panel solve is too slow for a
map of every charge.

The ion leaves along $+x$. The plates push it upward, and the flight leaves
through the \val{BareFace} edge of the world at $(\val{BareX}, \val{BareY})$,
at $t = \val{BareT}$, at an angle $\val{BareAng}$ degrees. The ordinate
$\val{BareY}$ stays below the detector, which begins at $\val{yDetLo}$. The
path only climbs, from the launch $\val{yLaunch}$ to $\val{BareYmax}$. The
sampled clearance of the plates is $\val{BareDmin}$, the energy drifts by
$\val{BareWdev}$, and the run takes $\val{BareSteps}$ steps.

The reference charge is positive and sits below the launch, so it repels the
ion upward, and it rearranges the charge on the plates through
\eqref{eq:collocation}. That extra lift is the difference between the two
curves in Figure~\ref{fig:board}: the reference meets the detector still
rising, with clearance $\val{RefDmin}$.

\section{What the game computes}\label{sec:engine}

The engine solves the same boundary condition (\physref{2.7}). The surfaces are
cut into flat triangles with constant density, graded toward the edges, and the
potential of a triangle is the closed integral named there. Only the upper half
is meshed, and the mirror is applied exactly. The linear system is a dense LU
factorization. Preview flights use panels of half a cell, and verification
repeats them at a finer grade. Pictures use a coarser mesh and are not the
flight.

This document collocates the same condition, equation \eqref{eq:collocation},
on equal rectangles of edge $\val{panelFine}$ with the integral
\eqref{eq:rectangle}. The factorization is again a dense LU. The end states
move by $\val{errMesh}$ when that edge doubles, and the differences in
\S\ref{sec:comparison} are the same size. The two meshes are not the same
discretization, so the agreement asked of them is the tolerance
$\val{tolFlight}$, ten times that estimate.

The ion's image force is omitted on both sides. The bound sampled here reaches
$\val{RefImage}$, under the engine's limit $10^{-10}$. The obstacle is the box
plus the contact shell $\val{contact}$. Because the panel field is a gradient,
$W$ in \eqref{eq:flight-energy} drifts only by $\val{RefWdev}$ on the reference
flight.

The flight itself, in the engine, is the adaptive eighth-order Runge--Kutta
method of \physref{5} and \physref{6}, verified as in \physref{7}. Once the
field is fixed, a verified end state is good to about $10^{-10}$. Our step
estimate $\val{errStep}$ is of that order; the larger difference is the
panels. Radiation is integrated and not applied to the momentum
(\xsecref{foundations}{sec:radiation}, \physref{8}). The fraction
$\val{RefRadFrac}$ is our integral of our field.

\section{Further reading}\label{sec:reading}

The motion, the energy, the Runge--Kutta step, the Hermite location of an
event and the Simpson rule are in the Foundations
(\xsecref{foundations}{sec:numerics}, \xeqref{foundations}{eq:energy},
\xsecref{foundations}{sec:radiation}). The uniform gap \eqref{eq:uniform} and
the rectangle collocation \eqref{eq:collocation} are derived above; later plate
levels cite them. The engine's electrodes are \physref{2.7}.

\levelfileappendix

\end{document}
